\section{Mode Preference and Free-Text Feedback}
\label{sec:eval-freeform}

The per-instrument data, so far, have shown a consistent Form-leading
pattern on usability (Section~\ref{sec:eval-sus}) and workload
(Section~\ref{sec:eval-rtlx}), with a partial inversion on perceived
transparency, trust, and perceived confidence
(Section~\ref{sec:eval-ai-support}). This section analyzes how the three
patterns combine when participants are forced to commit to a single
overall preference. After all three interaction modes had been
completed, participants provided two final inputs: an overall
ranking of the three modes from most to least preferred, and an
open-ended summary comment on the study experience. In addition,
after each individual mode, an optional per-mode comment was
collected.

\subsubsection{Mode Preference Ranking}

The ranking offers the cleanest test of which mode participants
actually preferred once they had used all three.
Table~\ref{tab:mode-preference} and Figure~\ref{fig:ranking-distribution} summarize
the post-experiment ranking.
Participants assigned each of the three modes a rank from
$1$ (most preferred) to $3$ (least preferred). ``Rank-1 votes''
counts the participants who placed the mode first.

\begin{table}[htbp]
\centering
\caption{Post-experiment mode preference ranking ($n = 69$).
``Rank-1 votes'' counts the number of participants who placed the
mode first; ``Share'' is the corresponding percentage; ``Mean rank''
averages the assigned $1$--$3$ rank across all $69$ participants.}
\label{tab:mode-preference}
\begin{tabular}{@{}l r r r@{}}
\toprule
Mode & Rank-1 votes & Share & Mean rank \\
\midrule
Form & \textbf{37} & \textbf{53.6\,\%} & \textbf{1.78} \\
Chat & 20 & 29.0\,\% & 2.07 \\
Ask  & 12 & 17.4\,\% & 2.14 \\
\bottomrule
\end{tabular}
\end{table}

Form received the plurality of first-place votes ($37$ of $69$, $53.6\,\%$),
followed by Chat ($20$, $29.0\,\%$) and Ask ($12$, $17.4\,\%$),
with mean ranks following the same ordering ($1.78$, $2.07$, $2.14$).
The preference for Form is statistically significant, with a
medium effect size.\footnote{$\chi^2(2) = 14.17$, $p < 0.001$,
Cram\'{e}r's $V = 0.32$. Full statistics in
Appendix~\ref{app:inferential}.}

\begin{figure}[htbp]
\centering
\includegraphics[width=0.85\linewidth]{ranking_distribution.png}
\caption{Distribution of all three rank positions per mode
($n = 69$).}
\label{fig:ranking-distribution}
\end{figure}

\subsubsection{Free-Text Comments}

The free-text comments contextualize the ranking by surfacing the
reasons participants gave for their preferences. Each mode collected
an optional per-mode open-ended comment, and after all three modes
participants wrote an overall summary comment. A total of $88$
comments were recorded across $44$ unique participants. Comments are
written in German with a small minority in English; the
paragraphs that follow summarize them thematically in the order
Form, Ask, Chat, and close with the across-modes overall
comments.

\paragraph{Form.} Comments on Form cluster around speed, clarity,
and control. Recurring observations describe the mode as the fastest
of the three, the easiest to obtain an overview of, and the most
``integrated'' into the document creation workflow, as the
checkbox interface makes the selection space visible at a glance.
Several participants emphasize that Form ``takes the most work off
you'' (\emph{``nimmt dir am meisten Arbeit ab''}) by pre-populating
the option set, while a smaller subset notes the opposite concern:
that the absence of an interactive LLM commentary makes it unclear on
what basis the model produces its suggestions, and that participants
would not feel safe relying on Form output without an additional
verification step. Several comments describe specific UX issues:
long text fields rendered as a single narrow line on mobile, the
inability to provide the LLM with additional context for its
suggestions, and the impression that the LLM ``guesses'' without
scenario-specific input.

\paragraph{Ask.} Comments on Ask are the most polarized of the
three modes. A consistent theme is that Ask is the most cumbersome
mode in terms of workflow: the requirement to manually copy the LLM's
suggested text into the form field (cf.\ Section~\ref{sec:ropagen}) is
repeatedly described as redundant work, particularly as
participants reported that the model often produced suggestions in a
free-text style that did not match the structured form fields.
Several participants observed that
the lack of a persistent context across fields forced them to
re-enter the same information repeatedly. Conversely, explanatory
text generated by Ask is consistently described as the most
informative of the three modes, supporting the high transparency
rating visible in item~3 of the AI Support questionnaire. A frequent
suggestion is to combine Ask's explanations with Form's structured
field layout.

\paragraph{Chat.} Comments on Chat are dominated by two themes:
comfort and opacity. On the comfort side, several participants
describe Chat as the most intuitive of the three modes, especially
as the LLM takes over the work of mapping the conversation onto
the structured ROPA fields (cf.\ Section~\ref{sec:ropagen}), leaving the
participant only with the task of describing the processing scenario. On the opacity side, an
equally large group of participants reports uncertainty about what is
actually being written into the document — whether the entire chat
history is used for generation, whether each individual message must
be confirmed, and what the document looks like during the
conversation. Specific defects mentioned include the LLM occasionally
switching from German to English mid-response, occasional emission of
raw \texttt{DATA\_UPDATE} tokens, and a slow ``streaming'' rendering
style that participants experienced as artificially throttled.

\paragraph{Overall comments} The overall comments
recapitulate the per-mode pattern. A clear majority describes Form as
the fastest and most ergonomic mode, particularly for participants
who feel they ``do not know enough'' to interrogate the model
meaningfully. Chat is typically placed second and motivated by the
impression that ``the AI takes care of it,'' which several
participants explicitly note increases their \emph{subjective} sense
of safety even though they cannot verify the output. Ask is most
often placed third, with the manual copy-paste step cited as the
principal reason. A recurring constructive suggestion across
participants is a hybrid mode that combines the visible option set of
Form with the conversational, context-aware behavior of Chat or
Ask. Several participants explicitly add a caveat that they are
novices in the legal domain and that their preference ordering might
shift if they had prior GDPR knowledge.

The ranking and comment data align with the \textit{SUS} and \textit{R-TLX}
picture. Form is the most-chosen mode by a clear margin; Chat
slots in second on raw votes; Ask is most often placed third,
with free-text comments citing the manual copy-paste step between
explanation and form field as the principal friction. Across the
comments the most consistent constructive suggestion is for a
hybrid mode that combines Form's visible option set with Chat's
or Ask's conversational behavior.
